% --- Archon physics formalization source begin ---
% archon:physics
% archon:covers IPhO2026Problems/problem_ipho_2026_e1_c6.lean
% archon:source-report reports/ipho_2026/problem_ipho_2026_e1_c6.source.json
% archon:problem-id ipho_2026_e1
% archon:part-id E1-C6
% archon:previous-part ipho_2026_e1_c5
% archon:previous-part-policy derive_inline_from_problem_only_material

\chapter{Ipho Physics problem ipho\_2026\_e1\_c6}
\label{ch:IPhO2026Problems_problem_ipho_2026_e1_c6}

\paragraph{Problem source.}
\#\# Physical scenario

Water in the inner and outer cylinders exchanges heat radially through an
acrylic cylindrical wall.  Record T\_IC and T\_OC versus time.  The heat-flow
model is dQ/dt = (T\_OC - T\_IC)/R\_Th.  For radial Fourier conduction,
dQ/dt = -lambda*A*dT/dr.  Ignore apparatus heat capacity where instructed and
use the dimensions in Figure 17.

\#\# Current subquestion E1-C6

Determine the effective wall thermal resistance R\_Th from the C5 graph.

\paragraph{Shared problem context.}
Water in the inner and outer cylinders exchanges heat radially through an
acrylic cylindrical wall.  Record T\_IC and T\_OC versus time.  The heat-flow
model is dQ/dt = (T\_OC - T\_IC)/R\_Th.  For radial Fourier conduction,
dQ/dt = -lambda*A*dT/dr.  Ignore apparatus heat capacity where instructed and
use the dimensions in Figure 17.

\paragraph{Current subquestion.}
Determine the effective wall thermal resistance R\_Th from the C5 graph.

\paragraph{Figure/image path.}
ipho\_2026\_source/image/E1\_page-13.png

\paragraph{Reusable previous-part conclusions.}
\begin{itemize}
\item Source E1-C5. Question: Graph the finite-difference rate (T\_IC,j-T\_IC,j-1)/(t\_j-t\_j-1) against the corresponding average T\_OC-T\_IC. Policy: derive\_inline\_from\_problem\_only\_material
\end{itemize}

\paragraph{Formalization target.}
create a compiling Lean file with sorry bodies at `IPhO2026Problems/problem\_ipho\_2026\_e1\_c6.lean`.
The Lean declarations must preserve the physical quantities, dimensions or dimensional roles, figure labels, governing-law hypotheses, and final relation expressed by this problem.
Use Mathlib/Physlib names found through LeanExplore where available. If a domain API is missing, introduce faithful local abstractions rather than scalar placeholder aliases.
This is an answer-blind solve-phase task. Derive a candidate result only from the problem-side evidence above; do not assume a target value, select a tolerance around a desired value, or encode a desired result into a candidate domain.
For numerical questions, define the raw end-to-end quantity and apply an explicit source-derived rounding rule. For identification questions, characterize or prove uniqueness before recording the derived candidate.

\begin{theorem}[Ipho Physics formalization target]
\label{thm:physics:ipho_2026_e1_c6:target}
\leanok
The assigned autoformalize agent should translate this IPhO physics problem into Lean declarations in the covered file, with theorem and lemma proof bodies written as `by sorry`.
\end{theorem}
\begin{proof}
\leanok
Fully proved in \texttt{IPhO2026Problems/problem\_ipho\_2026\_e1\_c6.lean}
(namespace \texttt{IPhO2026.E1C6}): the two law fields of
\texttt{IntervalHeatExchange} (equation (4) in finite-interval form and the
IC-water calorimetry $\Delta Q = m_{IC} \cdot c_w \cdot \Delta T_{IC}$)
eliminate $\Delta Q$ to the combined interval equation
$m \cdot c_w \cdot \Delta T_{IC} = (\bar T_{OC} - \bar T_{IC})/R_{Th} \cdot \Delta t$
(\texttt{ic\_water\_heating\_eq\_wall\_flow}); dividing by the positive time
step and by $m \cdot c_w \cdot R_{Th}$ gives the C5 proportionality with
physical slope $1/(m \cdot c_w \cdot R_{Th})$
(\texttt{finiteDifferenceRate\_eq\_physical\_slope}), and cancelling the
nonzero abscissa identifies the graph readout $k = 1/(m \cdot c_w \cdot R_{Th})$
(\texttt{c5GraphSlope\_eq\_physical}). Inverting with the positive IC water
mass (\texttt{icWaterMass\_pos}) yields the target
$R_{Th} = 1/(m_{IC} \cdot c_w \cdot k)$
(\texttt{thermalResistance\_eq\_candidate}). The heating-branch certificates
\texttt{heatReceived\_pos}, \texttt{finiteDifferenceRate\_pos},
\texttt{c5GraphSlope\_pos} and \texttt{determined\_resistance\_pos} are closed
by \texttt{positivity} along the same laws. All seven former \texttt{sorry}
bodies are now kernel-checked proofs; no new axioms beyond
\texttt{propext}, \texttt{Classical.choice}, \texttt{Quot.sound}.
\end{proof}
% --- Archon physics formalization source end ---
